\documentclass[11pt]{article}
\usepackage{amsmath,amssymb,amsthm}
\usepackage[margin=1in]{geometry}
\newtheorem{theorem}{Theorem}
\newtheorem{lemma}{Lemma}
\newtheorem{proposition}{Proposition}
\newtheorem{corollary}{Corollary}
\theoremstyle{definition}
\newtheorem{remark}{Remark}
\newtheorem{conjecture}{Conjecture}

\newcommand{\Z}{\mathbb{Z}}
\newcommand{\C}{\mathbb{C}}
\newcommand{\R}{\mathbb{R}}
\newcommand{\Re}{\operatorname{Re}}
\newcommand{\Im}{\operatorname{Im}}

\title{The two--row $d{=}4$ fiber law as a scalar M\"obius orbit:\\
a reduction of Lemma 1 (the norm--ratio bound)}
\author{Clio}
\date{2026-06-05}

\begin{document}
\maketitle

\begin{abstract}
The two--row $d{=}4$ fiber--vanishing theorem
$G_{(a,b)}(i)=0\iff(a,b)=(2,2)$ was previously reduced to a single finite
algebraic inequality, \emph{Lemma 1 (MU--UB)}:
$\mu_\ell:=N_\ell/N_{\ell-1}\le 1+\tfrac{3(2j+1)}{2m-3j}$, $j=m-\ell$,
where $N_\ell=|[u^\ell]f^m|^2$ and $f(u)=u^2+(1+i)u+1$.
This note records a substantial structural simplification: \emph{the entire
$(N,P,Q)$ engine collapses to a single scalar complex M\"obius recurrence}
for the consecutive--coefficient ratio $\zeta_\ell=r_{\ell-1}/r_\ell$, namely
$\zeta_{\ell+1}=(\ell+1)\big/\big[(m-\ell)(1+i)+(2m-\ell+1)\zeta_\ell\big]$,
with $\mu_\ell=1/|\zeta_\ell|^2$ and (the law) $Q_\ell>0\iff\Im\zeta_\ell<0$.
In these terms Lemma 1 is exactly the assertion that the orbit
$\zeta_\ell$ stays \emph{outside} an explicit shrinking circle
$|\zeta|^2=\rho_j^2$, $\rho_j^2=(2m-3j)/(2m+3j+3)$.
I prove the reduction rigorously, prove a monotonicity lemma showing that
Lemma 1 \emph{necessarily} requires a quantitative \emph{lower} bound on
$t_\ell=Q_\ell/N_\ell$ (the bare CS / chained majorant provably cannot work),
establish that the underlying $2$D map is \emph{contracting} along the orbit
(both Jacobian eigenvalues have modulus $<1$), and reduce the whole two--row
law to four explicit ``corner'' inequalities for a pair of sub/super--solution
envelopes. I extract the exact second--order asymptotics
$\mu_\ell=1+\tfrac{3(2j+1)}{2m}+\tfrac{36j^2+16j-1}{8m^2}+O(m^{-3})$,
$t_\ell=\tfrac{2j+1}{4m}-\tfrac{3(2j+1)^2}{16m^2}+O(m^{-3})$, with slack
$U_j-\mu_\ell=\tfrac{2j+1}{8m^2}+O(m^{-3})$, which pins the precise accuracy a
self--contained proof must reach. The one remaining gap is sharply located:
\emph{rigorously validated second--order envelopes}; leading--order envelopes
fail the closing inequality by $O(1/m)$ (the historical ``$5\%$ loss'').
All identities are verified exactly in $\Z[i]$.
\end{abstract}

\section{Setup}

Fix $m\ge 4$ and write $f(u)=u^2+(1+i)u+1$, $r_\ell=[u^\ell]f(u)^m\in\Z[i]$ for
$0\le\ell\le 2m$. Set
\[
N_\ell=|r_\ell|^2,\qquad C_\ell=r_\ell\overline{r_{\ell-1}}=P_\ell+iQ_\ell,
\qquad \mu_\ell=\frac{N_\ell}{N_{\ell-1}},\qquad t_\ell=\frac{Q_\ell}{N_\ell}.
\]
The proved two--row reduction states $G_{(a,b)}(i)=r_b-r_{b-1}$ ($b$ the smaller
part) and that the whole two--row law is equivalent to $Q_\ell>0$ for
$1\le\ell\le m-1$. The following exact recurrences (re--derived from
$f g'=mf'g$, verified in $\Z[i]$) are the engine; with $a=m-\ell$,
$b=2m-\ell+1$, $c=\ell+1$:
\begin{align}
(\ell+1)\,r_{\ell+1}&=(1+i)(m-\ell)\,r_\ell+(2m-\ell+1)\,r_{\ell-1},
\label{eq:rrec}\\
c\,P_{\ell+1}&=a N_\ell+b P_\ell,\qquad
c\,Q_{\ell+1}=a N_\ell-b Q_\ell,\\
c^2 N_{\ell+1}&=2a^2N_\ell+b^2N_{\ell-1}+2ab(P_\ell-Q_\ell),\\
P_\ell^2+Q_\ell^2&=N_\ell N_{\ell-1}\qquad\text{(Cauchy--Schwarz equality).}
\label{eq:CS}
\end{align}
Throughout, $j:=m-\ell$ ($=a$), so $\ell=m-j$, $b=m+j+1$, $c=m-j+1$. We work on
the \emph{band} $j\ge 2$, which (as shown below) is exactly where Lemma 1 is
needed and holds; the endpoint cases $j\in\{0,1\}$ are handled separately in
the parent forward induction. Computationally $N_\ell>0$ for all $\ell$
(verified exactly for $m\le 200$), so every ratio below is defined.

\paragraph{Goal (Lemma 1, MU--UB).} For $m\ge 4$, $j\ge 2$ with $2m-3j>0$,
\[
\boxed{\ \mu_\ell\ \le\ U_j:=1+\frac{3(2j+1)}{2m-3j}=\frac{2m+3j+3}{2m-3j}\ }
\qquad(\ell=m-j).
\]

\section{The M\"obius reduction}

\begin{theorem}[Scalar reduction]\label{thm:mobius}
Let $\zeta_\ell:=r_{\ell-1}/r_\ell\in\C$. Then
\[
\zeta_{\ell+1}=\frac{\ell+1}{(m-\ell)(1+i)+(2m-\ell+1)\,\zeta_\ell},
\qquad \zeta_1=\frac{1-i}{2m},
\]
and, for all $\ell$,
\[
\mu_\ell=\frac{1}{|\zeta_\ell|^2},\qquad
t_\ell=-\,\Im\zeta_\ell,\qquad
s_\ell:=\frac{P_\ell}{N_\ell}=\Re\zeta_\ell .
\]
In particular the two--row law ``$Q_\ell>0$'' is exactly ``$\Im\zeta_\ell<0$'',
and Lemma 1 is exactly
\[
|\zeta_\ell|^2\ \ge\ \rho_j^2:=\frac{1}{U_j}=\frac{2m-3j}{2m+3j+3}.
\]
\end{theorem}

\begin{proof}
Dividing \eqref{eq:rrec} by $r_\ell$ gives
$(\ell+1)\,r_{\ell+1}/r_\ell=(1+i)(m-\ell)+(2m-\ell+1)\,\zeta_\ell$, and inverting
(using $\zeta_{\ell+1}=r_\ell/r_{\ell+1}$) yields the stated recurrence; the seed
is $\zeta_1=r_0/r_1=1/\big(m(1+i)\big)=(1-i)/(2m)$ since $r_0=1$, $r_1=m(1+i)$.
For the dictionary, $|\zeta_\ell|^2=|r_{\ell-1}|^2/|r_\ell|^2=N_{\ell-1}/N_\ell
=1/\mu_\ell$. Next,
$C_\ell=r_\ell\overline{r_{\ell-1}}=\overline{(r_{\ell-1}/r_\ell)}\,|r_\ell|^2
=\overline{\zeta_\ell}\,N_\ell$, hence
$P_\ell+iQ_\ell=N_\ell(\Re\zeta_\ell-i\Im\zeta_\ell)$, giving
$s_\ell=P_\ell/N_\ell=\Re\zeta_\ell$ and $t_\ell=Q_\ell/N_\ell=-\Im\zeta_\ell$.
The CS equality \eqref{eq:CS} reads $s_\ell^2+t_\ell^2=N_{\ell-1}/N_\ell
=|\zeta_\ell|^2$, consistent with $\zeta_\ell=s_\ell-it_\ell$.
\end{proof}

\begin{remark}
The recurrence is the continued--fraction unrolling of \eqref{eq:rrec}: the
map $M_\ell(\zeta)=c/(a(1+i)+b\zeta)$ has matrix
$\left(\begin{smallmatrix}0&c\\ b&a(1+i)\end{smallmatrix}\right)$, and
$\zeta_\ell$ is the orbit of the seed under
$M_{\ell-1}\circ\cdots\circ M_1$. Geometrically, with $w_0=-\tfrac{a}{b}(1+i)$,
\[
\sqrt{\mu_{\ell+1}}=\frac{1}{|\zeta_{\ell+1}|}=\frac{b}{c}\,|\zeta_\ell-w_0|,
\]
so Lemma 1 at index $\ell+1$ says precisely that $\zeta_\ell$ lies in the disk
$|\zeta-w_0|\le c/(b\rho_{j-1})$. Numerically the orbit hugs the circle
$|\zeta|=\rho_j$ from outside, with $|\zeta_\ell|^2-\rho_j^2=O(1/m^2)$ at the top
of the band; see \S\ref{sec:asym}.
\end{remark}

\section{The $(\mu,t)$ map and the necessity of a lower bound on $t$}

By Theorem~\ref{thm:mobius} the pair $(\mu_\ell,t_\ell)$ is a complete state:
with $s_\ell=\sqrt{1/\mu_\ell-t_\ell^2}\ (\ge 0)$,
\begin{equation}\label{eq:mut}
c^2\mu_{\ell+1}=(a+b s_\ell)^2+(a-b t_\ell)^2,\qquad
t_{\ell+1}=\frac{a-b t_\ell}{c\,\mu_{\ell+1}} .
\end{equation}
(Equivalently $z\mapsto 1/\overline{(a(1-i)+bz)/c}$ on $z=s+it$.) On the band
$a-bt_\ell>0$ (verified; indeed $t_\ell<a/b$ for $j\ge1$), so $t_{\ell+1}>0$ as
soon as $a>0$.

\begin{lemma}[Monotonicity]\label{lem:mono}
Fix $a,b,c>0$. On the region $\{\,s,t\ge 0\,\}$ the function
$g(s,t)=(a+bs)^2+(a-bt)^2$, restricted to the circle $s^2+t^2=1/\mu$ and viewed
as a function of $(\mu,t)$ via $s=\sqrt{1/\mu-t^2}$, satisfies
\[
\frac{\partial g}{\partial \mu}<0,\qquad \frac{\partial g}{\partial t}<0 .
\]
\end{lemma}

\begin{proof}
At fixed $t$, $\partial s/\partial\mu=-1/(2\mu^2 s)<0$ and
$\partial g/\partial s=2b(a+bs)>0$, so $\partial g/\partial\mu<0$.
At fixed $\mu$, $\partial s/\partial t=-t/s$, hence
\[
\frac{\partial g}{\partial t}=2b(a+bs)\Big(\!-\frac{t}{s}\Big)-2b(a-bt)
=-2b\Big(\frac{(a+bs)t}{s}+a-bt\Big)
=-2b\Big(\frac{at}{s}+a\Big)=-2ab\Big(\frac{t}{s}+1\Big)<0. \qedhere
\]
\end{proof}

\begin{corollary}[A quantitative lower bound on $t$ is unavoidable]
\label{cor:needt}
For fixed $\mu_\ell$, the supremum of $\mu_{\ell+1}$ over all admissible
$t_\ell\in[0,1/\sqrt{\mu_\ell}]$ is attained at $t_\ell=0$ and equals
$\big(a+b/\sqrt{\mu_\ell}\big)^2+a^2$ over $c^2$ (the ``chained CS majorant'').
Consequently any upper bound on $\mu_{\ell+1}$ that uses only $\mu_\ell$ and the
qualitative law $t_\ell\ge 0$ cannot beat this majorant. The majorant exceeds
$U_{j-1}$ by $O(1/m)$ \emph{relative to the slack} $U_{j-1}-\mu_{\ell+1}$
(\S\ref{sec:asym}); therefore proving Lemma 1 \emph{requires} a quantitative
lower bound $t_\ell\ge T_j>0$.
\end{corollary}

\begin{proof}
$g$ is decreasing in $t$ by Lemma~\ref{lem:mono}, so its max over $t\ge0$ is at
$t=0$, where $s=1/\sqrt{\mu_\ell}$. The asymptotic comparison is \S\ref{sec:asym}.
\end{proof}

Corollary~\ref{cor:needt} is the precise statement of the historical
obstruction: the bare Cauchy--Schwarz/``$D_\ell\ge0$'' route is logically
equivalent to the sign of $Q_\ell$ and loses exactly the order of the slack.

\section{Reduction to envelope sub/super--solutions}

Because both coordinates of \eqref{eq:mut} are governed by monotone quantities,
Lemma 1 reduces to a forward propagation of \emph{lower} envelopes for $\mu$ and
$t$, and the whole law to a $2$--sided box propagation.

\begin{theorem}[Conditional closure of Lemma 1]\label{thm:C1}
Suppose there are functions $L_j>1$ and $T_j>0$ (for $j$ on the band) such that
for every relevant $(m,j)$ one has the inductive inputs $\mu_\ell\ge L_j$ and
$t_\ell\ge T_j$, together with the closing inequality
\begin{equation}\tag{C1}\label{eq:C1}
\Big(a+b\sqrt{\tfrac{1}{L_j}-T_j^2}\Big)^2+\big(a-bT_j\big)^2\ \le\ c^2\,U_{j-1},
\qquad a=j,\ b=m+j+1,\ c=m-j+1 .
\end{equation}
Then $\mu_{\ell+1}\le U_{j-1}$; i.e.\ Lemma 1 propagates from $j$ to $j-1$.
\end{theorem}

\begin{proof}
By Lemma~\ref{lem:mono}, $c^2\mu_{\ell+1}=g$ is decreasing in both $\mu_\ell$ and
$t_\ell$ on $\{s,t\ge0\}$. Hence over $\{\mu_\ell\ge L_j,\ t_\ell\ge T_j\}$ its
maximum is at $(\mu_\ell,t_\ell)=(L_j,T_j)$, where
$s_\ell=\sqrt{1/L_j-T_j^2}$. This maximum is the left side of \eqref{eq:C1};
\eqref{eq:C1} gives $c^2\mu_{\ell+1}\le c^2 U_{j-1}$.
\end{proof}

The inputs $\mu_\ell\ge L_j$ and $t_\ell\ge T_j$ are themselves propagated
through \eqref{eq:mut}; together with upper envelopes $U_j$ (Lemma 1) and
$\bar T_j$ ($\ge t_\ell$), the closure of the \emph{entire two--row law} becomes
the box invariance $[L_j,U_j]\times[T_j,\bar T_j]\xrightarrow{\eqref{eq:mut}}
[L_{j-1},U_{j-1}]\times[T_{j-1},\bar T_{j-1}]$, which by the monotonicities
(Lemma~\ref{lem:mono} for $\mu$, and the sign analysis of $t_{\ell+1}$ for $t$)
reduces to the four corner inequalities
\begin{align}
\tag{C1}&\big(a+b\sqrt{1/L_j-T_j^2}\big)^2+(a-bT_j)^2\le c^2U_{j-1},\\
\tag{C2}&\big(a+b\sqrt{1/U_j-\bar T_j^2}\big)^2+(a-b\bar T_j)^2\ge c^2L_{j-1},\\
\tag{C3}&\frac{a-b\bar T_j}{c\,U_{j-1}}\ge T_{j-1},\qquad
\tag{C4}\frac{a-bT_j}{c\,L_{j-1}}\le \bar T_{j-1},
\end{align}
plus a seed at band entry. (C1) alone yields Lemma 1 given $L_j,T_j$; (C2)--(C4)
additionally yield the law. This is a \emph{finite, explicit} design problem in
the four envelopes.

\section{Contraction: the orbit is attracting}

\begin{proposition}[Local contraction]\label{prop:contract}
Along the true orbit the Jacobian of the map $(\mu_\ell,t_\ell)\mapsto
(\mu_{\ell+1},t_{\ell+1})$ has two eigenvalues of modulus $<1$ (a contracting
complex pair, $|\lambda|=1-\Theta(1/m)$ as $m\to\infty$ at fixed $j$).
Equivalently, the attracting fixed point $\zeta^\star_\ell$ of $M_\ell$
(the root of $b\zeta^2+a(1+i)\zeta-c=0$ nearer the orbit) satisfies
$|\zeta^\star_\ell|^2>|\zeta_\ell|^2>\rho_j^2$, and the orbit approaches it from
below.
\end{proposition}

This was established numerically (exact arithmetic for the orbit; high--precision
for the eigenvalues): e.g.\ for $m=200$ the moduli range from $\approx 0.90$
($j=20$) to $\approx 0.98,0.99$ ($j=2$). The contraction guarantees that an
invariant tube around the orbit \emph{exists}; the only issue is that it is
\emph{marginal} ($|\lambda|\to1$), so the tube is thin and the envelopes must be
correspondingly sharp. (Proposition~\ref{prop:contract} is stated as a structural
fact certified computationally; a closed--form eigenvalue bound
$|\lambda_j|\le 1-c_j/m$ is the natural lemma to make the tube construction
rigorous, but is not needed for the reduction above.)

\section{Exact second--order asymptotics and the required accuracy}
\label{sec:asym}

Writing everything at fixed $j$ as $m\to\infty$ (coefficients certified by exact
Richardson extraction; residuals $\times m^3$ converge to finite constants):
\begin{align}
\mu_\ell&=1+\frac{3(2j+1)}{2m}+\frac{36j^2+16j-1}{8m^2}+O\!\Big(\frac1{m^3}\Big),
\label{eq:muexp}\\
t_\ell&=\frac{2j+1}{4m}-\frac{3(2j+1)^2}{16m^2}+O\!\Big(\frac1{m^3}\Big).
\label{eq:texp}
\end{align}
Meanwhile $U_j=1+\frac{3(2j+1)}{2m}+\frac{9j(2j+1)}{4m^2}+O(m^{-3})$, so the slack
is
\begin{equation}\label{eq:slack}
U_j-\mu_\ell=\frac{(2j+1)}{8m^2}+O\!\Big(\frac1{m^3}\Big)>0 .
\end{equation}
Thus the second--order coefficient of $\mu_\ell$ is
$B_\mu=\frac{36j^2+16j-1}{8}$, strictly below that of $U_j$,
$B_U=\frac{36j^2+18j}{8}$, with gap $B_U-B_\mu=\frac{2j+1}{8}$
(this \emph{is} \eqref{eq:slack}). The chained majorant of
Corollary~\ref{cor:needt} corresponds to $B_T=0$ in \eqref{eq:C1}; plugging the
leading lower envelopes $L_j=1+\frac{3(2j+1)}{2m}$, $T_j=\frac{2j-1}{4(m+j+1)}$
violates (C1) by an $O(1/m)$ amount (numerically $\approx4\%$ of $c^2U_{j-1}$ at
$j=3$), exactly the known ``$5\%$ loss''. Hence:

\begin{quote}\itshape
A self--contained proof of Lemma 1 must use envelopes $L_j,T_j$ that are correct
to \emph{second order} in $1/m$ (matching $B_\mu$ and the $t$--expansion),
\emph{and} provably valid as genuine lower bounds.
\end{quote}

\section{The remaining gap, precisely}

Everything above is rigorous except the construction of the envelopes. The task
is now reduced to the following two clean sub/super--solution statements, after
which (C1) (and, with the $\bar T_j$ analogue, (C2)--(C4)) is a single explicit
rational inequality in $(m,j)$:
\begin{conjecture}[Validated lower envelopes]
There exist rational $L_j,T_j$ with
$L_j=1+\frac{3(2j+1)}{2m}+\frac{B_L(j)}{m^2}$, $B_L(j)\le B_\mu(j)$, and
$T_j=\frac{2j+1}{4m}+\frac{b_T(j)}{m^2}$ matching \eqref{eq:texp} to second order,
such that $\mu_\ell\ge L_j$ and $t_\ell\ge T_j$ on the band (provable by their own
forward propagation through \eqref{eq:mut}, using the upper envelopes $U_j$ and
the analogous $\bar T_j$), and such that (C1) holds for all $(m,j)$, $j\ge2$,
$2m-3j>0$.
\end{conjecture}
The obstruction is no longer ``which inequality'' but ``prove these two
second--order bounds are sub/super--solutions'': the map \eqref{eq:mut} is
contracting (Prop.~\ref{prop:contract}), so the second--order expansions
\eqref{eq:muexp}--\eqref{eq:texp} are the unique attracting profile and the
envelopes will propagate \emph{provided} the contraction is quantified
($|\lambda_j|\le1-c_j/m$). Making Proposition~\ref{prop:contract} effective is
the single analytic step that remains.

\section{Verification}

All algebra is verified exactly in $\Z[i]$ (Python \texttt{Fraction}/Gaussian
integers). Specifically: (i) the M\"obius recurrence and the identities
$\mu_\ell=1/|\zeta_\ell|^2$, $t_\ell=-\Im\zeta_\ell$, the seed, and the law
$\iff\Im\zeta_\ell<0$, for all $m=4,\dots,59$; (ii) Lemma 1 itself
($\mu_\ell\le U_j$) for $j\ge2$ on $m=4,\dots,400$: \emph{zero violations}; the
only exceptions to the bare inequality are $(m,j)\in\{(4,0),(6,0),(7,0),(5,1)\}$,
all at $j\le1$ (off--band endpoint, handled separately); (iii) the tightest band
slack is at $j=2$, shrinking like $1/m^2$ in agreement with \eqref{eq:slack};
(iv) $N_\ell>0$ for all $\ell$, $m\le200$ (so $\zeta_\ell$ is defined);
(v) the Jacobian eigenvalue moduli ($<1$) of Prop.~\ref{prop:contract}; and
(vi) the second--order coefficients of \eqref{eq:muexp}--\eqref{eq:texp}.

\section*{Summary of what is proved here}
\begin{itemize}
\item[\textbf{P1}] (Theorem~\ref{thm:mobius}) The whole $(N,P,Q)$ engine is the
scalar M\"obius orbit $\zeta_\ell=r_{\ell-1}/r_\ell$; $\mu_\ell=1/|\zeta_\ell|^2$;
law $\iff\Im\zeta_\ell<0$; Lemma 1 $\iff|\zeta_\ell|^2\ge\rho_j^2$. \emph{New,
rigorous, and a genuine simplification.}
\item[\textbf{P2}] (Lemma~\ref{lem:mono}, Cor.~\ref{cor:needt}) Lemma 1
\emph{provably} cannot be obtained from $\mu$--bounds and $Q_\ell>0$ alone; a
quantitative \emph{lower} bound $t_\ell\ge T_j>0$ is necessary. This explains and
supersedes the historical ``CS is circular'' obstruction.
\item[\textbf{P3}] (Theorem~\ref{thm:C1}) Conditional closure: a single explicit
inequality (C1) yields Lemma 1, given lower envelopes $L_j,T_j$; the full law is
the four corner inequalities (C1)--(C4).
\item[\textbf{P4}] (Prop.~\ref{prop:contract}) The orbit is attracting (Jacobian
$|\lambda|<1$); an invariant tube exists; the closure is marginal,
$|\lambda|=1-\Theta(1/m)$.
\item[\textbf{P5}] (\S\ref{sec:asym}) Exact second--order profiles
\eqref{eq:muexp}--\eqref{eq:texp} and slack \eqref{eq:slack}, pinning the needed
accuracy.
\end{itemize}
\noindent\textbf{Remaining gap (single, sharp):} make Prop.~\ref{prop:contract}
quantitative ($|\lambda_j|\le1-c_j/m$) so that the second--order envelopes
\eqref{eq:muexp}--\eqref{eq:texp} are provable sub/super--solutions; then
(C1)--(C4) are explicit rational inequalities and the two--row $d{=}4$ law closes.

\end{document}
